% ECONOMICS_PROBLEM: {"id": "O31-633", "title": "Direction and novelty", "jel_code": "O31", "source_id": "OP-0419"}
% !TeX program = xelatex
\documentclass[11pt,a4paper]{article}
\usepackage[margin=23mm]{geometry}
\usepackage{fontspec}
\setmainfont{Latin Modern Roman}
\usepackage{amsmath,amssymb,mathtools}
\usepackage{xcolor,enumitem,booktabs,longtable,array,xurl,hyperref}
\definecolor{muted}{HTML}{53636C}
\hypersetup{colorlinks=true,urlcolor=blue,linkcolor=blue}
\setlength{\parindent}{0pt}
\setlength{\parskip}{5pt}
\newcommand{\R}{\mathbb R}
\newcommand{\E}{\mathbb E}
\newcommand{\N}{\mathbb N}
\newcommand{\ind}{\mathbf 1}
\newcommand{\Var}{\operatorname{Var}}
\newcommand{\Cov}{\operatorname{Cov}}
\newcommand{\supp}{\operatorname{supp}}
\DeclareMathOperator*{\argmin}{arg\,min}
\DeclareMathOperator*{\argmax}{arg\,max}
\newcommand{\entrymeta}[3]{{\sffamily\small\color{muted}\textbf{\texttt{#1}}\quad|\quad #2\par #3\par}}
\newcommand{\Problemref}[1]{\texttt{#1}}
\newcommand{\JELref}[2]{\texttt{#2} (JEL \texttt{#1})}
\begin{document}
\section*{Direction and novelty}
\textbf{Problem ID:} \texttt{O31-633}\quad \textbf{JEL:} \texttt{O31}\par
% BEGIN ECONOMICS STATEMENT
\label{problem:OP-0419}
{\sffamily\small\color{muted}Original subject group: Innovation and AI economics\par}
\label{definitions:problem:30007701}
\textbf{Audit classification:} Background; proposed extension. The source is an automation/patenting study, not direct provenance for a randomized AI novelty theorem. Novelty scores and the experiment are explicitly editorial.
\subsection*{Definitions and precise research target}
Study research projects begun by firms in a specified technology sector during a defined enrollment year. Let \(A_i\in\{0,1\}\) denote randomized access to an AI discovery tool, with the same baseline research budget in each arm. For each project, define \(V_i(a)\geq0\) as externally validated scientific or commercial value over three years and \(R_i(a)\in[0,1]\) as novelty relative to a frozen pre-intervention knowledge corpus, judged using a preregistered blinded procedure. Failed projects have \(V_i(a)=0\). Define novel discovery value \(Y_i(a)=V_i(a)R_i(a)\), total-value effect \(\tau_V=E[V_i(1)-V_i(0)]\), and novelty-weighted effect \(\tau_Y=E[Y_i(1)-Y_i(0)]\). Also record projects' disciplinary and application categories before outcomes to measure changes in research direction without redefining categories after treatment. Estimate these quantities for eligible projects, accounting for connected research-team spillovers through cluster assignment. Patent counts or textual distance alone are insufficient substitutes for validated discovery value. Does AI increase \(\tau_Y\), or mainly expand valuable refinements with low measured novelty, and which project-selection changes mediate the difference? This specified experiment extends existing patent evidence; neither the novelty score nor a universal direction-of-innovation claim is supplied as an established open theorem.

\textbf{Definitions, assumptions and mathematical qualifications.} An eligible project has the same inclusion rule and baseline resource allowance under both policies, with AI/model access precisely specified. Novelty \(R_i(a)\) is a fixed scoring functional comparing a validated result to a frozen baseline corpus and blinded review protocol; it is not an intrinsic property inferred from word distance. Require finite \(E[V_i(a)]\) and define failed-project novelty convention, with \(V_iR_i=0\) regardless of that convention. For a fixed technology-category partition, report the joint distribution of output category, value and novelty. A project-selection mediator is itself treatment affected and needs a distinct causal intervention to identify mediation.
\subsection*{Variants and assumption changes}
\begin{enumerate}
\item \textbf{Editorial, existing projects.} Randomize AI access after project selection and estimate gross and novelty-weighted validated value effects for this fixed cohort.
\item \textbf{Editorial, direction choice.} Randomize access before firms allocate an identical research budget among a fixed category set; measure reallocations and market-level novel value with cluster spillovers, without conditioning on the chosen project category.
\end{enumerate}
\subsection*{References and source audit}
\textbf{Checked source.} NBER Working Paper 34240 (September 2025), indexed primary abstract. The cited work studies automation and innovation shifts; it does not supply this AI novelty experiment. Primary metadata and abstract indexed; direct record returned an access error. \url{https://www.nber.org/papers/w34240}.
\begin{enumerate}\raggedright
\item\hypertarget{definitions:source:30007701:1}{} Lin William Cong; Yao Lu; Hanqing Shi; Wu Zhu. 2025. \emph{Automation-Induced Innovation Shift}. NBER Working Paper 34240 (September 2025), indexed primary abstract.
\url{https://www.nber.org/papers/w34240}.
DOI: \url{https://doi.org/10.3386/w34240}.
\end{enumerate}

\subsection*{Common conventions: Source mathematical conventions}

These conventions apply to each entry unless explicitly replaced.

\textbf{Models and identification.} A primitive model \(m\) specifies all probability laws, feasible choices, information, equilibrium and measurement rules used by its target. The admissible class \(\mathcal M\) is fixed independently of outcomes. If \(P_O(m)\) is its observable probability law and \(\tau(m)\) the target, its identified set at observable law \(P\) is
\[
 \mathcal I_\tau(P)=\{\tau(m):m\in\mathcal M,\ P_O(m)=P\}.
\]
Point identification means that this set is a singleton. An empty set signals incompatibility of data and assumptions. Coordinate bounds need not describe the joint set. Bounds are sharp only if every retained target is attainable or arbitrarily approachable by compatible models. Finite bounds require explicit restrictions on unobserved quantities; unrestricted confounding, hidden wealth, extrapolation or arbitrary equilibrium selection can yield uninformative sets.

\textbf{Causal interventions.} A potential outcome \(Y_i(p)\) is the outcome under a complete policy regime \(p\), including timing, financing, information and market spillovers. The target population and units are fixed before treatment. With interference, use the complete assignment vector or a justified exposure mapping; individual no-interference notation alone cannot identify an economy-wide intervention. A difference of mean potential outcomes does not require the cross-world joint distribution. Conditioning on post-treatment migration, survival, enrollment or employment changes the estimand and needs separate justification.

\textbf{Statistical statements.} An estimator, confidence set, forecast or test specifies a sampling law, model class and loss/coverage criterion. Uniform \((1-\alpha)\)-coverage means \(\inf_{m\in\mathcal M}\Pr_m\{\tau(m)\in C_n\}\geq1-\alpha\), or an explicitly asymptotic version. The whole parameter space trivially has coverage, so informativeness requires a stated width, risk or power objective. A forecasting improvement is relative to a named benchmark, information set and evaluation population.

\textbf{Equilibrium and optimization.} An equilibrium comprises feasible plans, optimality, beliefs where needed, budget constraints and all market/resource clearing. The primitives or a stated benchmark define each correspondence \(\mathcal E(p;m)\). Nonemptiness is assumed or proved, never inferred just from compact policy sets. A selected-equilibrium target uses a declared selection rule; a robust target quantifies over all permitted equilibria. Use suprema or infima unless attainment is established. An \(\varepsilon\)-optimal policy is feasible and within \(\varepsilon\) of the finite optimum value. Report an empty feasible set separately.

\textbf{Welfare and accounting.} Normative utility, interpersonal normalization, social weights, numeraire, discounting and the use of public receipts are primitives. Behavioral utility need not coincide with the welfare criterion. Household welfare includes ownership income and rents once; adding profits again to compensating variation generally double-counts them. A monetary equivalent is defined by an explicit transfer and price convention, with monotonicity and range conditions. Average effects, mechanism contributions, transfers and welfare losses are different targets. Infinite sums require convergence and dynamic feasible plans require terminal or no-Ponzi conditions.

\textbf{Source versions.} A working-paper date, revision date and journal publication year are distinguished. A DOI for a journal version is not automatically the DOI of an earlier manuscript. Locators refer to the stated version and printed pages where specified. A metadata-only check does not verify a claim about a full-text conclusion. Every entry's source subsection narrows the provenance claim accordingly.
% END ECONOMICS STATEMENT
\end{document}
